\chapter{Weierstrass preparation}

This section is dedicated to the classical theorem "Weierstrass preparation".

%% Need to convert everything into BGR notation and work... then hope it can be translate to
%% use for Gouvea's work

\begin{theorem}\label{WP}
    Let $c$ be a positive real number and $f (x) = \sum a_n x^n$ be a power series with coefficients
    in $K$ such that $ \lvert a_n \rvert c^n \to 0$ as $n \to \infty$. Let $N$ be the number defined
    by the conditions
    \[
    \lvert a_N \rvert c^N = \max \lvert a_n \rvert c^n = \|f(x)\|_c \quad \text{and} \quad
    \lvert a_n \rvert c^n < \lvert a_N \rvert c^N \text{ for all } n > N.
    \]
    Then there exist a polynomial
    \[
    g(x) = b_0 + b_1x + \cdots + b_N x^N
    \]
    of degree $N$ and with coefficients in $K$, and a power series
    \[
    h(x) = 1 + d_1 x + d_2 x^2 + \cdots
    \]
    with coefficients in $K$, satisfying:
    \begin{itemize}
        \item $f(x) = g(x) h(x)$;
        \item $\lvert b_N \rvert c^N = \max \lvert b_n \rvert c^n$, that is
        $\| g (x) \|_c = \lvert b_N \rvert c^N$;
        \item $h(x) \in \mathcal{A}_c$, that is it is a restricted power series over $K$ for
        parameter $c$;
        \item $\lvert d_n \rvert c^n < 1$ for all $n \geq 1$, so that $\| h(x) - 1 \|_c < 1$; and
        \item $\| f(x) - g(x) \|_c < \| f(x)\|_c.$
    \end{itemize}
    In particular, $h(x)$ has no zeroes in $\overline{B}(0,c)$.
\end{theorem}

To prove this we will need to give a string of lemmas that will all come together.

We begin by showing that the Tate algebra with parameter $c$ is complete with respect to the Gauss
norm of parameter $c$, \cite[Lemma~7.2.7]{Gouvea}.

\begin{lemma}
    $\mathcal{A}_c$ is complete with respect to the Gauss norm $\|\cdot \|_c$.
\end{lemma}

\begin{proof}
    We need to show a Cauchy sequence in $\mathcal{A}_c$ converges with respect to the norm
    $\|\cdot\|_c$. Therefore, consider a sequence of power series
    \[
    f_i(x) = a_{i,0} + a_{i,1}x + a_{i,2}x^2 + \cdots.
    \]
    Saying this sequence is Cauchy means
    \[
    \forall \; \varepsilon>0, \; \exists \; M \text{ such that } \forall \; i,j > M
    \text{ we have } \|f_i (x) - f_j (x) \|_c < \varepsilon.
    \]
    Using the definition of the Gauss norm, this means that
    \[
    \max_{n} |a_{i,n}-a_{j,n} | c^n < \varepsilon \quad \forall \; i,j > M,
    \]
    and certainly,
    \[
    |a_{i,n} - a_{j,n}| < \varepsilon c^{-n}
    \]
    for each $n$ whenever $i,j > M.$ That is each of the sequences $(a_{i,n})_{i \in \N}$ forms a
    Cauchy sequence. Since $K$ is complete, this means they are convergent.
    Now consider
    \[
    g(x) = a_0 + a_1 x + a_2 x^2 + \cdots
    \]
    where $a_n = \lim_{i \to \infty} a_{i,n}$. We want to show this is the limit of the $f_n$ and
    is also in $\mathcal{A}_c$.

    For the first part, we know that if $i,j > M$ we have
    \[
    |a_{i,n} - a_{j,n} | c^n < \varepsilon
    \]
    and letting $j \to \infty$ it follows that if $i > M$
    \[
    | a_{i,n} - a_n| c^n \leq \varepsilon.
    \]
    Therefore, we have $\forall \; i > M$
    \[
    \| f_i (x) - g(x) \|_c = \max_n |a_{i,n} - a_n| c^n \leq \varepsilon
    \]
    and so $f_i(x) \to g(x)$ with respect to $\| \cdot \|_c.$

    Finally, since $f_ i \in \mathcal{A}_c$, there exists $N$ such that for all
    $n > N$, $|a_{i,n} | c^n \leq \varepsilon.$ Thus for all $n > N$
    \[
    | a_n |c^n = | a_n - a_{i,n} + a_{i,n}|c^n  \leq \max \{|a_n - a_{i,n}|c^n,|a_{i,n}|c^n\} \leq \varepsilon
    \]
    and therefore $|a_n|c^n \to 0$, which by definition means $g(x) \in \mathcal{A}_c$.
    With this we conclude the proof.
\end{proof}

Next we show that polynomials are dense in the Tate algebra, \cite[Lemma~7.2.8]{Gouvea}.
\begin{lemma}\label{PolyDense}
    The space of polynomials $K[x]$ is dense in $\mathcal{A}_c$.
\end{lemma}

\begin{proof}
    We need to show that any power series is the limit of a sequence of polynomials. Towards this
    let
    \[
    f(x) = a_0 + a_1x+a_2x^2 + \cdots
    \]
    be a power series in $\mathcal{A}_c$.
    Consider the sequence of truncations of $f(x)$,
    \begin{align*}
        f_0(x) &= a_0 \\
        f_1(x) &= a_0 + a_1 x \\
        f_2(x) &= a_0 + a_1 x + a_2 x^2 \\
        & \; \; \vdots \\
        f_k(x) &= a_0 + a_1 x + \cdots a_k x^k
    \end{align*}
    we want to show $f$ is the limit of this sequence. Now
    \[
    \| f(x) - f_k (x) \|_c = \max_{n>k} |a_n|c^n,
    \]
    which tends to zero as $k \to \infty$ since $\lim_{n \to \infty} |a_n|c^n = 0$. Thus
    $f_k (x) \to f(x)$ with respect to the Gauss norm and we conclude.
\end{proof}

Finally, we give a generalisation of \cite[Lemma~7.2.9]{Gouvea}, which is an analogue of the
Euclidean algoritheorem for Tate algebras.

\begin{lemma} \label{euclidalgoritheorem}
    Let $f(x) \in \mathcal{A}_c$, and let
    \[
    g(x) = b_0 + b_1 x + \cdots + b_N x^N
    \]
    be a polynomial with coefficients in $K$ satisfying
    \[
    \lvert b_N \rvert  c^N= \max_i (\lvert b_i \rvert c^i).
    \]
    Then there exist a power series $q(x) \in \mathcal{A}_c$ and a polynomial $r(x) \in K[x]$,
    of degree less that $N$, such that
    \[
    f(x) = g(x) q(x) + r(x)
    \]
    where $q(x)$ and $r(x)$ satisfy
    \[
    \|f(x) \|_c \geq \|g(x)\|_c \| q(x)\|_c \quad \text{and} \quad \|f(x) \|_c \geq \| r(x) \|.
    \]
    \end{lemma}

\begin{proof}
    Need to write up.
\end{proof}

With these theorems we are now in a position to prove the $p$-adic Weierstrass preparation theorem.

\begin{proof}[Proof of Theorem~\ref{WP}]
    The idea is to start with an approximate factorisation, and then improve it. If we prove by
    induction this is always possible, it will produce convergent sequences of polynomials and power
    series which in the limit gives the factorisation we want.

    \textbf{Base case:}
    To start the induction we need to find a $\delta < 1$ and an initial pair $g_1 (x)$ and
    $h_1 (x)$. Towards this consider
    \[
    g_1 (x) = \sum^{N}_{i = 0} a_i x^i \quad h_1(x) = 1.
    \]
    Then, since the assumption is that $\|f(x)\|_c = |a_N |c^N$ and the terms of higher degree are
    smaller, we have
    \[
    \| f(x) - \sum_{i=0}^N a_i x^i\|_c < \|f(x)\|_c
    \]
    and let $\delta$ be a measure of how much smaller:
    \[
    \|f(x) - \sum_{i=0}^N a_i x^i \|_c = \delta \|f(x) \|_c,
    \]
    clearly $0 < \delta < 1$. Then we can easily check that all our properties are satisfied.

    \textbf{Inductive step:}
    Let $\delta$ be a fixed real number, $0 < \delta < 1$. Suppose at stage $n$ we have a polynomial
    $g_n (x)$ and a power series $h_n (x)$ such that
    \begin{enumerate}
        \item $\|f(x) - g_n (x) h_n(x) \|_c \leq \delta^n \|f(x)\|_c$;
        \item $\deg g_n (x) = N$;
        \item $\|g_n(x) \|_c = |b_{n,N} | c^N$;
        \item $h_n (x) \in \mathcal{A}_c$;
        \item $\| h_n(x) - 1 \|_c \leq \delta^n$; and
        \item $\|f (x) - g_n (x) \|_c < \delta \|f(x)\|_c.$
    \end{enumerate}

    Firstly, $\|f(x) - g_n (x) \|_c \leq \delta \|f(x) \|$ implies $ \| f(x) \|_c = \| g(x)_n \|_c$,
    using the fact $\| \cdot \|_c$ is non-archimedean and that $0 < \delta < 1$.

    By Lemma \ref{euclidalgoritheorem} we can find a polynomial $r(x)$ of degree less than $N$ and a
    power series $q(x) \in \mathcal{A}_c$ such that
    \[
    f(x) - g_n (x) h_n (x) = q(x) g_n(x) + r(x)
    \]
    and
    \begin{align*}
    \| q(x) \|_c & \leq \frac{\| f(x) - g_n (x) h_n (x) \|_c}{\| g_n (x) \|_c} \leq
    \frac{\delta^n \|f_n(x)\|_c}{\|f(x) \|_c} = \delta^n \leq \delta\\
    \| r(x) \|_c &\leq \| f(x) - g_n (x) h_n (x) \|_c \leq
    \delta^n \| f(x) \|_c \leq \delta \| f(x) \|_c.
    \end{align*}
    Now consider
    \[
    g_{n+1} = g_n (x) + r(x) \quad \text{and} \quad h_{n+1}(x) = h_n (x) + q(x).
    \]
    We claim these are better approximations and satisfy the wanted properties.

    Firstly, since $\deg r(x) < N$ we must have $\deg g_{n+1} (x) = N = \deg g_n (x)$ and
    $g_{n+1}$ will also have leading term $b_{n,N}$.
    Moreover, since $\mathcal{A}_c$ forms a ring, $h_{n+1} \in \mathcal{A}_c$ (polynomials are
    trivially restricted power series).

    Next, we have
    \begin{align*}
        \| f(x) - g_{n+1}(x) \|_c &= \| f(x) - g_n (x) - r(x) \|_c \\
        & \leq \max \{ \|f(x) - g_n(x) \|_c, \| r(x) \|_c \} \\
        & \leq \delta \| f(x) \|_c
    \end{align*}
    which means by the same argument as for $g_n$
    \[
    \|g_{n+1} (x)\|_c = \| f(x) \|_n = \| g_n ( x) \|_c = |b_{n,N}| c^N = | b_{n+1,N} | c^N.
    \]
    We also have
    \begin{align*}
        \| h_{n+1} (x) - 1 \|_c & = \| h_n (x) - 1 + q(x) \|_c \\
        & \leq \max \{ \| h_n (x) - 1 \|_c , \| q(x) \|_c \} \\
        & \leq \delta .
    \end{align*}

    Now we want to show it indeed gives a better approximation:
    \begin{align*}
        f(x) &- g_{n+1} (x) h_{n+1}(x) = f(x) - (g_n (x) + r(x) ) (h_n(x) + q(x)) \\
        &= f(x) - g_n (x) h_n (x) - q(x) g_n(x) - r(x) h_n (x) - r(x) q(x) \\
        & = r(x) - r(x) h_n (x) - r(x) q(x) \\
        & = r(x) ( 1 - h_n (x) - q(x) )
    \end{align*}
    and therefore
    \begin{align*}
    \| f(x) - g_{n+1} (x) h_{n+1} (x) \|_c &= \| r(x) \|_c \|(1 - h_n(x)) - q(x) \|_c \\
    & \leq \delta \|f(x) \|_c \max\{ \|(1 - h_n (x)) \|_c, \|q(x)\|_c \} \\
    & \leq \delta^{n+1} \|f(x) \|_c.
    \end{align*}
    Since, $0 < \delta < 1$ this implies we do indeed have a better approximation.

    Therefore $g_{n+1}$ and $h_{n+1}$ satisfies all our properties with $\delta^n$ replaced by
    $\delta^{n+1}$.

    To conclude, we want to show these functions actually form the sequence we want.
    Firstly, they are Cauchy as
    \[
    \| q(x) \|_c \leq \delta^n \implies \|h_n (x) - h_{n+1}  (x) \|_c \leq \delta^n
    \]
    and
    \[
    \| r(x) \|_c \leq \delta^i \|f(x) \|_c \implies
    \|g_n (x) - g_{n+1} (x) \|_c \leq \delta^i \| f(x) \|.
    \]
    Therefore, as $\mathcal{A}_c$ is complete, we have they converge to a $g(x)$ and $h(x)$. It is
    then immediate via taking limits that these have the wanted properties.
\end{proof}
